%
% 2-liouville.tex
%
% (c) 2026 Prof Dr Andreas Müller
%
\section{Das Prinzip von Liouville
\label{buch:intalg:section:liouville}}
\kopfrechts{Das Prinzip von Liouville}%
Die in Kapitel~\ref{chapter:ratint} durchgerechneten Beispiele 
uggerieren, dass eine Stammfunktion immer nur eine Linearkombination
von Logarithmen der Teilnenner der Partialbruchzerlegung ist.
Das Prinzip von Liouville zeigt, dass dies ganz allgemein gilt.

%
% Das Prinzip
%
\subsection{Das Prinzip 
\label{buch:intalg:liouville:subsection:prinzip}}

\begin{satz}
Sei $F$ ein Funktionenkörper mit Konstantenkörper $K$ und $f\in F$.
Wenn $f$ eine Stammfunktion $g$ in einer elementaren Erweiterung
$G\supset F$ hat, die den gleichen Konstantenkörper wie $F$ hat, dann
gibt es $v_0,v_1,\dots,v_n\in F$ und $c_1,\dots,c_k\in K$ mit
\begin{equation}
f
=
Dv_0
+
\sum_{i=1}^n c_i\frac{Dv_i}{v_i},
\end{equation}
d.~h.~
\[
\int f(z)\,dz
=
v_0
+ 
\sum_{i=1}^n c_i \log(v_i).
\]
$G$ muss also alle Logarithmen $\log(v_i)$, $i=1,\dots,n$, enthalten.
\end{satz}

Wie kann man überhaupt eine solche Aussage beweisen?
Wenn $g$ die Stammfunktion von $f$ ist, dann muss die Ableitung
von $g$ wieder $f$ ergeben, also $Dg=f$.
Wenn also $g$ in einer Erweiterung $F(\vartheta_1,\dots,\vartheta_n)$
von $F$ liegt, dann kann man versuchen, einen Ausdruck für $g$ in den
Elementen $\vartheta_k$ abzuleiten.
Dabei muss ein Ausdruck entstehen, welcher mit $f$ übereinstimmt.
Alle Terme mit den Elementen $\vartheta_k$ müssen sich daher wegheben.

Der allgemeine Fall einer Erweiterung der Form
$F(\vartheta_1,\dots,\vartheta_n)$ führt zu ziemlich komplizierten
Ausdrücken.
Daher beginnen wir damit einer Erweiterung $F(\vartheta)$ um ein
einziges $\vartheta$.
In diesem Fall muss $g$ die Form einer rationalen Funktion
\[
g
=
\frac{a(\vartheta)}{b(\vartheta)} 
=
\frac{
a_0 + a_1\vartheta + \dots a_n\vartheta^n
}{
b_0 + b_1\vartheta + \dots b_m\vartheta^m
},
\]
wobei $a(\vartheta)\in F[\vartheta]$ und $b(\vartheta)\in F[\vartheta]$
Polynome mit Koeffizienten in $F$ sind.
In diesem Fall lässt sich die Ableitung von $g$ leichter 
berechnen.
Zu diesem Zweck berechnen wir im
Abschnitt~\ref{buch:intalt:liouville:subsection:polynome}
zunächst die Ableitung von Polynomen von $\vartheta$ für die
verschiedenen Arten von Erweiterungen.
Diese Resultate verwenden wir dann im
Abschnitt~\ref{buch:intalt:liouville:subsection:spezialfaelle},
um das Prinzip von Liouville für eine Erweiterung der Form $F(\vartheta)$
zu beweisen.
In Abschnitt~\ref{buch:intalt:liouville:subsection:allgemein}
wird dann der allgemeine Fall bewiesen.

%
% Polynome von Erweiterungen
%
\subsection{Polynome von Erweiterungen
\label{buch:intalg:liouville:subsection:polynome}}

%
% Beweis des Prinzips von Liouville in Spezialfällen
%
\subsection{Beweis des Prinzips von Liouville in Spezialfällen
\label{buch:intalg:liouville:subsection:spezialfaelle}}

%
% Beweis im allgemeinen Fall
%
\subsection{Beweis im allgemeinen Fall
\label{buch:intalg:liouville:subsection:allgemein}}



